
\chapter{Криптография в~платежах}\label{ch:crypto}
Платёжная криптография отвечает на~три инженерных вопроса. Что
произойдёт, если из~терминала утечёт ключ. Почему нельзя хранить
PIN в~базе и~сверять его программно. Как карта и~эмитент доверяют
друг другу, если секреты нельзя раздать всем участникам цепочки.

\section{Симметричная криптография: 3DES и AES}\label{sec:crypto-symmetric}

В~симметричной криптографии один и~тот же ключ работает на~шифрование
и~расшифрование, поэтому отправитель и~получатель заранее договариваются
о~секрете и~хранят его в~тайне. Платёжная индустрия исторически
строилась на~симметричных алгоритмах: они быстрые, компактные
и~хорошо подходят к~задачам, где обе стороны уже находятся внутри
одной доверенной инфраструктуры.

\subsection*{3DES: рабочая лошадка унаследованных систем}

3DES, он же Triple DES или TDEA (Triple Data Encryption Algorithm),~---
тройное применение DES, разработанного в~1970-х годах.
Одиночный DES с~56-битным ключом давно непригоден для практической
защиты; тройное применение с~двумя или тремя ключами заметно
поднимает стойкость и~сохраняет совместимость с~парком устройств,
рассчитанных на~DES.

В~платёжной инфраструктуре 3DES встречается повсюду. На~нём считают
CVV магнитной полосы (раздел~\ref{sec:card-data-cvv}), генерируют
криптограмму ARQC (Authorization Request Cryptogram, см.~\autoref{sec:emv-cryptogram}) на~чипе, шифруют
PIN-блок при онлайн-авторизации, вычисляют MAC (Message
Authentication Code, имитовставка в~отечественной терминологии)
для защиты целостности ISO~8583-сообщений
в~канале между эквайером и~сетью. За~десятилетия 3DES проник
вглубь терминалов, HSM (Hardware Security Module, аппаратный модуль безопасности; подробно в~\autoref{sec:crypto-hsm}) и~процессинговых систем. Его замена
стала долгой и~дорогой миграцией. NIST (National Institute of Standards and Technology) отозвал SP~800--67 Rev.~2 (SP~--- Special Publication) с~1~января 2024~года; после этого TDEA по~линии NIST больше не~считается утверждённым блочным шифром и~сохраняется лишь
в~ограниченных унаследованных сценариях: расшифрование, снятие
ключевой обёртки и~проверка MAC уже защищённых данных~\cite{nist-sp800-67r2,nist-sp800-67-withdrawal}.

\subsection*{AES: современная замена}

AES (Advanced Encryption Standard)~--- блочный шифр, принятый NIST
в~2001~году на~замену DES~\cite{nist-fips-197}.
Он работает с~128-битными блоками и~поддерживает длину ключа
128, 192 или 256~бит. AES быстрее 3DES, требует меньше
вычислительных ресурсов и~даёт значительно больший запас
стойкости для новых систем.

AES постепенно вытесняет 3DES в~платёжной инфраструктуре.
Новые терминалы и~модули HSM поддерживают AES штатно, но~миграцию
тормозит обратная совместимость с~унаследованным парком устройств.

\subsection*{MAC: защита целостности сообщений}

MAC (Message Authentication Code)~--- значение, вычисляемое
по~содержимому сообщения на~разделяемом секретном ключе. Получатель
по~нему проверяет, что сообщение не~изменилось в~пути и~пришло
от~стороны, владеющей ключом. В~ISO~8583 MAC обычно передают
в~DE~64 или DE~128 и~используют для защиты авторизационного
сообщения на~участке между терминалом, эквайером и~сетью.

В~платёжной практике чаще всего встречаются два алгоритма MAC:
CBC-MAC поверх 3DES или AES в~режиме CBC (Cipher Block Chaining, сцепление блоков шифротекста) и~более современный
CMAC (Cipher-based MAC, NIST SP~800-38B) на~AES. Выбор определяется тем,
что уже умеет терминальный парк и~что поддерживает сторона эмитента.

\section{Асимметричная криптография: RSA и ECC}\label{sec:crypto-asymmetric}

В~симметричной криптографии один ключ на~двоих. Асимметричная
использует пару: открытый ключ распространяется свободно, закрытый
хранится в~тайне. Данные, зашифрованные открытым ключом,
расшифровываются только закрытым; подпись, созданная закрытым,
проверяется открытым. Асимметрия устанавливает доверие между
сторонами, которые заранее не~обменивались секретами~\cite{emvco-book2}.

\subsection*{RSA в~EMV}

RSA остаётся самым узнаваемым асимметричным алгоритмом в~платёжной
индустрии. В~разделе~\ref{sec:emv-data-auth} показано, как RSA работает
для аутентификации данных на~чипе. Карточная сеть выпускает
сертификаты ключей, эмитент подписывает сертификат карты, терминал
проверяет цепочку доверия на~основе заранее загруженного открытого
ключа корневого центра.

С~момента появления EMV контактный чип опирался на~RSA, а~длины
ключей EMVCo пересматривает ежегодными бюллетенями
\enquote{RSA Key Lengths Assessment}: требования зависят от~роли
ключа, версии спецификации и~текущего миграционного
окна~\cite{emvco-ecc-faq,emvco-bulletins-archive}. Универсальная
рекомендация \enquote{RSA-NNNN бит} быстро устаревает.

\subsection*{ECC: короче ключи, легче для чипа}

ECC (Elliptic Curve Cryptography)~--- семейство алгоритмов на
эллиптических кривых. Главное преимущество перед RSA~--- сопоставимая
стойкость при существенно меньшем размере ключа: публичный ключ P-256
занимает 32~байта против 256~байт у~RSA-2048, сертификатная цепочка
короче. На~стороне карты эффект сильнее: подпись на~чипе быстрее
на~порядок и~ключевой материал занимает меньше EEPROM. На~стороне
терминала картина обратная: RSA verify с~малым публичным экспонентом
$e=65{\,}537$ обходится примерно 17~модульными умножениями и~работает
в~разы быстрее ECDSA P-256 verify. Но узкое место EMV-транзакции~---
подпись внутри карты в~коротком бесконтактном окне, а~не~verify
на~терминале, поэтому переход на~ECC ускоряет именно критичный
участок~\cite{emvco-ecc-faq}.

EMVCo официально добавил поддержку ECC в~спецификации контактного
чипа. Во~внутрироссийской карточной системе Мир ту же задачу цифровой
подписи решает ГОСТ~34.10--2018 на~эллиптических кривых~\cite{emvco-ecc-faq,gost-34-10}.
RSA не~исчезает сразу: переход идёт постепенно, потому что
терминалы, карты и~сертификационные цепочки должны одновременно
поддерживать старый и~новый набор алгоритмов.

\section{HSM}\label{sec:crypto-hsm}

\highlight{Любой криптографический алгоритм безопасен настолько, насколько
безопасен ключ}. Если ключ хранится файлом на~сервере, его уносят
вместе с~дампом памяти или резервной копией. Аппаратный модуль безопасности
HSM (Hardware Security Module)~--- физически защищённое устройство,
в~котором ключевой материал остаётся внутри корпуса, а~наружу
выходит только результат операции.

У~программной реализации ключ в~оперативной памяти достаётся
через дамп процесса, через \textit{cold-boot}-атаку
(извлечение ключей из~ОЗУ после быстрого охлаждения модулей и~аварийной перезагрузки)
или через доступ DMA с~вредоносного PCIe-устройства. Ни один
из~этих векторов не~требует взлома алгоритма; хватает физического
или привилегированного доступа к~машине.

HSM проектируется против физических атак: корпус обнаруживает
вскрытие, внутренние цепи защищают ключи от~считывания, побочные
каналы минимизированы. Физическую прочность удостоверяет формальная
сертификация: FIPS~140--3 (Federal Information Processing Standard, серия стандартов NIST) описывает четыре уровня
безопасности криптографических модулей, PCI~SSC (Payment Card Industry Security Standards Council)
отдельно поддерживает стандарт PCI~PTS~HSM (PIN Transaction Security Hardware Security Module
Modular Security Requirements) для платёжного контекста. Платёжному
HSM нужны все три: FIPS~140--3, PCI~PIN (PCI~PIN Security Requirements)
и~PCI~PTS~HSM~\cite{nist-fips-140-3,pci-pin-security,pci-pts-hsm}.

HSM обслуживает всю рутинную криптографию процессинга.
Через него проходит верификация криптограммы ARQC при
онлайн-авторизации: эмитент или его процессор передаёт данные
транзакции в~HSM, тот деривирует сессионный ключ из~мастер-ключа
эмитента, вычисляет ожидаемую криптограмму и~сравнивает с~полученной.
Через тот же модуль идут генерация ARPC (Authorization Response Cryptogram, ответная криптограмма эмитента; см.~\autoref{sec:emv-cryptogram}), верификация PIN с~расшифровкой
PIN-блока и~сверкой с~эталоном без выпуска PIN наружу,
перешифрование PIN-блока между зонами безопасности,
генерация и~проверка CVV/CVV2, подготовка ключей при персонализации
карт и~вычисление MAC для защиты ISO~8583-сообщений.

Маршрут ARQC от~терминала до~хоста эмитента~--- центральная ось
рис.~\ref{fig:crypto-hsm-flow}. HSM эквайера держит TPK/ZPK/CVK
под~LMK, транслирует PIN-блок между зонами (команды Thales payShield
\texttt{JC/JD}, \texttt{JE/JF}) и~проверяет CVV (\texttt{CW/CX});
зональный HSM перешифрует PIN-блок при смене ZPK (\texttt{CA/CB}).
HSM эмитента деривирует MK\_AC из~IMK-AC по~PAN+PSN, затем сессионный
SK\_AC по~ATC, верифицирует ARQC (\texttt{KQ}) и~генерирует
ARPC~=~MAC(SK\_AC, ARQC~XOR~ARC)~(\texttt{KR}). Формулы~--- из~Common
Session Key Derivation, Annex~A1.3 EMV Book~2~\cite{emvco-book2};
команды~--- из~Thales payShield Host Programmer's
Manual~\cite{thales-payshield-host-commands} (это отдельный документ
от~\enquote{Legacy Host Commands} с~другой схемой кодов AA/AE/CI/JS);
публичная сводка соответствий~--- у~EFTlab~\cite{eftlab-thales-hsm-commands}.

\begin{figure}[tbp]
  \centering
  \includefiguresvg[width=\linewidth]{assets/figures/ch10-cryptography/crypto-hsm-flow}
  \caption{HSM в~авторизационном потоке ARQC/ARPC: эквайер, зональный
  узел, эмитент.}%
  \label{fig:crypto-hsm-flow}
\end{figure}

В~реальной инфраструктуре HSM почти никогда не~стоит в~одиночку:
его включают в~отказоустойчивый пул, обслуживающий авторизацию,
персонализацию и~процессинг PIN. В~типичном потоке чиповой
авторизации HSM участвует и~на проверке ARQC, и~на формировании
ARPC, но точное число обращений зависит от~сети и~конфигурации.
В~процессинге пул HSM~--- жёсткое ограничение масштабирования
с~собственными правилами планирования ёмкости и~key
ceremony (см.~\ref{sec:pc-performance}).

\practicenote{%
  При отказе HSM авторизация останавливается. Без доступа
  к~мастер-ключу нельзя ни деривировать сессионный ключ, ни
  верифицировать ARQC. Стандартная схема~--- active-active кластер
  двух HSM с~синхронизированными ключами и~регулярной проверкой
  того, что оба устройства дают одинаковый результат на~контрольной
  операции.
}

\subsection*{HSM в~контуре Мир}\label{sec:crypto-hsm-mir}

Российский эмиссионный контур опирается на~ту же модель HSM,
но~под другими нормами и~с~гибридной криптографией. Требования
к~функциям модуля закреплены в~ФТТ Банка России №\,ФТ-56-3/35
\cite{cbr-ft-56-hsm}: поддержка ARQC/ARPC по~EMV~CSK, MAC
по~ГОСТ~Р~34.12/34.13, формат PIN-блока ISO~9564-1~Format~0.
Положение Банка России 719-П~\cite{cbr-719p} и~851-П~\cite{cbr-851p}
требуют применять СКЗИ, сертифицированные ФСБ; для~платёжного HSM
это класс КВ2~\cite{fsb-378}.

Сертифицированных моделей в~эмиссионном контуре сейчас три:
КриптоПро~HSM~2.0~R3 с~Платёжным модулем~\cite{cryptopro-hsm-payment},
ViPNet~HSM~PS~\cite{vipnet-hsm-ps} и~SPB~HSM~PS~\cite{spb-hsm-ps}.
Все три поддерживают одновременно и~ГОСТ~34.10/34.11/34.12/34.13, и~3DES/AES/RSA,
потому что карта Мир работает в~гибриде: онлайн-криптограммы ARQC
массово считаются по~3DES~CSK, а~ГОСТ применяется для~MAC канала
с~НСПК (на~ZAK) и~для~подписи клиринговых файлов с~УЦ~ПС~Мир.
Хост-API во~всех трёх вендорах эмулирует формат Thales payShield
для~совместимости с~устоявшимся процессинговым ПО~--- СмартВистой,
TranzWare, OpenWay.

Контур расходится с~международной схемой по~нескольким направлениям:
ФСБ~КВ2 вместо FIPS~140-3, MAC канала по~ГОСТ~34.13 вместо ISO~9797-1~alg.~3,
подпись клиринговых файлов на~ГОСТ~34.10--2018~+~Стрибог. Публичная
часть Программы безопасности ПС~Мир~\cite{nspk-mir-security}
конкретный список рекомендованных моделей HSM не~раскрывает.

\begin{figure}[tbp]
  \centering
  \includefiguresvg[width=\linewidth]{assets/figures/ch10-cryptography/crypto-hsm-mir}
  \caption{Контур Мир и~международная схема HSM:
  сертификация, ключи, MAC, подпись клиринговых файлов.}\label{fig:crypto-hsm-mir}
\end{figure}

\section{Иерархия ключей}\label{sec:crypto-key-hierarchy}

Если бы эмитент использовал один и~тот же секрет для всех карт
и~всех транзакций, компрометация одного узла обернулась бы
системной аварией. Поэтому ключи в~платёжной криптографии
организуют в~иерархию~--- дерево, где каждый уровень порождает
ключи следующего через деривацию. Компрометация ключа нижнего
уровня не~должна ломать всю систему~\cite{nist-sp800-57}.

Без деривации эмитент хранил бы $N$ независимых случайных
ключей, по~одному на~каждую карту. При ста миллионах карт это
сто миллионов записей в~HSM; при смене оборудования или
восстановлении из~резервной копии всё это переносится заново.
Деривация снимает обе проблемы. Один IMK детерминированно
воспроизводит ключ любой карты, без дополнительного хранения
и~без передачи ключей между системами.

Типичная иерархия криптографии эмитента в~EMV:

\begin{enumerate}
  \item \textbf{Мастер-ключ (Issuer Master Key, IMK)}~--- секретный
    ключ эмитента, хранящийся исключительно в~HSM. Для каждого
    типа операции может существовать отдельный мастер-ключ.
    IMK-AC отвечает за криптограммы, IMK-SMI~--- за Secure Messaging
    Integrity (целостность команд эмитента к~чипу, issuer~scripts),
    IMK-SMC~--- за Secure Messaging Confidentiality (конфиденциальность
    тех же команд, например при удалённой смене PIN).
  \item \textbf{Ключ карты (Card Key)}~--- производный от~мастер-ключа
    через PAN и~порядковый номер карты PAN Sequence
    Number в~качестве входных данных для деривации. Ключ карты
    записывается в~чип при персонализации.
  \item \textbf{Сессионный ключ (Session Key)}~--- производный от~ключа
    карты через ATC (Application Transaction Counter)
    как диверсификатор. Сессионный ключ уникален для каждой
    транзакции и~используется для вычисления криптограммы.
\end{enumerate}

\begin{figure}[tbp]
  \centering
  \includefiguresvg[width=\linewidth]{assets/figures/ch10-cryptography/crypto-key-hierarchy}
  \caption{Дерево деривации EMV: IMK, ключ карты, сессионный ключ.}\label{fig:crypto-key-hierarchy}
\end{figure}

\highlight{Деривация работает как односторонняя функция: зная мастер-ключ,
вычисляют ключ любой карты; зная ключ карты, вычисляют сессионный
ключ для любого ATC}. Обратная операция невозможна, поэтому
компрометация сессионного ключа одной транзакции не~восстанавливает
ни ключ карты, ни тем более мастер-ключ. Компрометация ключа одной
карты не~раскрывает ключи других, поскольку деривация по~PAN даёт разные
результаты для разных номеров~\cite{emvco-book2}.

Компрометация мастер-ключа разрушает иерархию: все карты
эмитента оказываются скомпрометированы разом. Мастер-ключ никогда
не~покидает HSM в~открытом виде и~загружается при начальной настройке
через формальную процедуру Key Ceremony (раздел~\ref{sec:crypto-key-ceremony}).

\section{DUKPT}\label{sec:crypto-dukpt}

Иерархия ключей EMV (мастер-ключ $\to$ ключ карты $\to$ сессионный
ключ) хорошо работает для чиповых карт, где ключевой материал
защищён внутри чипа. POS-терминал, в~отличие от~чипа, стоит
на~прилавке магазина: его можно украсть, разобрать и~попытаться
извлечь ключ из~памяти. Если терминал хранит
мастер-ключ и~деривирует из~него сессионные ключи по~счётчику,
компрометация мастер-ключа раскроет всю иерархию.

DUKPT (Derived Unique Key Per Transaction)~--- схема управления
ключами, придуманная под эту задачу.
В~актуальной редакции ANSI X9.24--3 она описана как AES DUKPT, в~которой
Base Derivation Key (BDK) порождает уникальные начальные ключи
устройств, а~из них по~номеру транзакции выводятся рабочие ключи
отдельных операций~\cite{ansi-x9-24-3}. \highlight{Главное свойство схемы:
одноразовость транзакционных ключей}. Компрометация текущего ключа
не~даёт вычислить ключи прошлых операций.

Без DUKPT терминал запрашивал бы свежий ключ с~сервера на~каждую
транзакцию. Для терминала на~прилавке это неприемлемо:
POS должен работать при нестабильном или временно отсутствующем
соединении. Терминал с~DUKPT деривирует ключ автономно
из~локального счётчика, без обращения к~серверу в~момент
транзакции.

\begin{enumerate}
  \item Центральная система хранит BDK (Base Derivation Key)~--- корневой
    ключ, из~которого деривируются ключи для всех терминалов.
  \item При инициализации терминала из~BDK выводится уникальный
    начальный ключ конкретного устройства.
  \item Из этого начального ключа по~номеру транзакции терминал
    выводит рабочий ключ отдельной операции.
  \item После использования транзакционный ключ стирается и~больше
    не~применяется.
\end{enumerate}

Проследим деривацию на~конкретном примере AES DUKPT. Эквайер
хранит в~HSM 128-битный BDK. При инициализации терминала с~идентификатором
Initial Key ID (IKI, 8~байт) HSM вычисляет начальный ключ устройства
Initial Key (IK) одним блоком AES в~режиме ECB:
$$\text{IK} = \text{AES-ECB}_{\text{BDK}}(\text{DD}_{\text{init}}),$$
где $\text{DD}_{\text{init}}$~--- 16-байтовая структура \textit{derivation data}
со~служебными полями (\textit{version}, \textit{key usage} = \texttt{0x8001}
для Initial Key, \textit{algorithm} = AES-128, \textit{length} = 128) и~младшими
8~байтами, равными IKI~\cite{ansi-x9-24-3}. Терминал получает IK
и~удаляет любые следы BDK. С~этого момента BDK существует только в~HSM.

Дальше работает 32-битный transaction counter. Терминал держит
до~32~промежуточных ключей Intermediate Derivation Key (IDK)~---
по~одному на~каждый установленный бит счётчика. Рабочий ключ конкретной
операции выводится из~соответствующего IDK ещё одним шагом AES-ECB
с~другим derivation data:
$$\text{Working Key} = \text{AES-ECB}_{\text{IDK}}(\text{DD}_{\text{wk}}),$$
где \textit{key usage} в~$\text{DD}_{\text{wk}}$ равно \texttt{0x1000}
для PIN encryption key, \texttt{0x2000}~--- для MAC, \texttt{0x3000}~---
для data encryption~\cite{ansi-x9-24-3}. После транзакции счётчик
инкрементируется, набор IDK перестраивается под новые установленные
биты, а~использованный working key стирается. Серверная сторона
(хост эквайера или платёжная сеть), зная BDK и~полученный с~транзакцией
Key Serial Number $\text{KSN} = \text{IKI} \,\|\, \text{counter}$ (12~байт),
повторяет ту~же цепочку деривации и~получает тот~же working key,
не~храня промежуточного состояния.

Почему компрометация working key транзакции~$N$ не~раскрывает прошлое.
AES-ECB необратима без ключа, а~IDK прошлых позиций счётчика стираются
при инкременте и~больше не~хранятся на~устройстве. Атакующий
с~украденным терминалом получает текущие IDK, из~которых выводятся
только будущие working keys; прошлые ключи невосстановимы. После
переполнения 32-битного счётчика (порядка $4\cdot10^9$ транзакций)
терминал требует повторной инициализации с~новым~IK.

Реализация всех трёх шагов деривации укладывается в~один файл; функция
\texttt{make\_derivation\_data} строит 16-байтовый блок по~разметке из
рис.~\ref{fig:crypto-dukpt}, а~\texttt{aes\_ecb\_block}~--- единственный
криптографический примитив:

\codefile{Python}{samples/ch10-dukpt/dukpt.py}

Получение working key для конкретной транзакции~--- три вызова подряд,
без~хранения промежуточного состояния:

\begin{codeblock}{Python}
ik = derive_initial_key(bdk, iki)
idk = derive_idk(ik, iki, counter)
pin_key = derive_working_key(idk, iki, counter, KU_PIN_ENCRYPTION)
\end{codeblock}

На~официальных векторах из~Supplement to X9.24-3-2017~\cite{x9-24-3-test-vectors}
эти три вызова возвращают именно те значения, что подсвечены жёлтым
на~рис.~\ref{fig:crypto-dukpt}. Замена \texttt{KU\_PIN\_ENCRYPTION}
на~\texttt{KU\_MAC\_GENERATION} или \texttt{KU\_DATA\_ENCRYPTION} даёт
независимый working key того~же IDK для других операций (\textit{key separation}).

\begin{figure}[p]
  \centering
  {\renewcommand{\figuremaxheight}{.92\textheight}%
   \includefiguresvg[width=\linewidth]{assets/figures/ch10-cryptography/crypto-dukpt}}
  \caption{AES DUKPT по~ANSI X9.24-3:2017: деривация Initial Key из~BDK
    и~Working Key из~IDK.}\label{fig:crypto-dukpt}
\end{figure}

\section{Онлайн PIN-блок: форматы ISO 0, 1, 3, 4}\label{sec:crypto-pin-block}

Когда держатель вводит PIN на~клавиатуре терминала для
онлайн-авторизации, его нужно передать дальше так, чтобы
ни терминал, ни промежуточные узлы не~держали PIN в~открытом виде
дольше необходимого. Для этого PIN упаковывают в~PIN-блок:
структурированное значение длиной 8~байт для унаследованных форматов
или 16~байт для Format~4 на~основе AES, готовое к~передаче
по~сети без раскрытия PIN. Затем его шифруют ключом терминала
или зональным ключом и~отправляют в~DE~52 авторизационного
сообщения. Стандарт ISO~9564 определяет несколько форматов
PIN-блока; в~платёжной практике обычно говорят об~унаследованных
форматах и~о~более новом формате 4 на~основе AES~\cite{iso-9564}.
Форматы 0, 1, 3 и~4, разбираемые в~этом разделе, относятся
к~онлайн-сценарию: передача PIN-блока в~авторизационном
сообщении. Format~2 в~ISO~9564 предназначен для офлайн-проверки EMV и~здесь не~рассматривается.

\subsection*{Format 0 (ISO 9564 Format 0)}

Самый распространённый унаследованный формат. Он привязывает PIN
к~PAN: один и~тот же PIN даёт разные блоки на~разных картах,
что усложняет словарный перебор по~базе перехваченных шифротекстов.

Сборка показана на~рис.~\ref{fig:pin-block-format0} на~примере
PIN~=~\texttt{1234} и~тестового PAN Stripe
\texttt{4242424242424242}~\cite{stripe-testing}: PIN и~PAN упаковываются
в~два 8-байтовых поля (по 16 нибблов), которые XOR-ятся в~итоговый
блок. PIN-поле~--- формат-идентификатор \texttt{0}, длина PIN, цифры
PIN, заполнитель \texttt{F}. PAN-поле~--- четыре нибла нуля, далее
правые 12~цифр PAN без check digit. Результат XOR~--- 8-байтовый
блок, готовый к~шифрованию ключом PEK (PIN Encryption
Key)~\cite{iso-9564}.

\begin{figure}[tbp]
  \centering
  \includefiguresvg[width=\linewidth]{assets/figures/ch10-cryptography/pin-block-format0}
  \caption{Сборка PIN-блока Format~0: PIN-поле, PAN-поле, XOR двух
  полей.}\label{fig:pin-block-format0}
\end{figure}

Та же привязка работает против безопасности при утечке PAN.
Если злоумышленник знает PAN (а~он часто известен из~тех же логов),
задача восстановления PIN из~перехваченного PIN-блока упрощается.
При смене PAN, например при перевыпуске карты, PIN-блок меняется,
даже если сам PIN остался прежним.

\subsection*{Format 1 и~Format 3}

Format~1 не~привязан к~PAN: он предназначен для случаев, где
PIN шифруется без зависимости от~номера карты. В~карточных потоках
встречается заметно реже, чем Format~0 и~Format~3.

Format~3 остаётся унаследованным, но ослабляет повторяемость блока:
вместо постоянного заполнителя использует переменные значения,
и~одинаковые PIN и~PAN дают разные выходы в~разных транзакциях.
Это отсекает часть статистических атак на~перехваченные
блоки~\cite{iso-9564}.

\subsection*{Format 4 (AES)}

Format~4~--- подход на~основе AES, задуманный как замена 3DES-логике
для новых систем~\cite{iso-9564}.

\importantnote{%
  Процессинг должен одновременно поддерживать AES-контур и~legacy; сроки PCI PIN Security пересматриваются~\cite{pci-format4-suspension,pci-format4-supplement}~--- не~закладывайте даты в~код.
}
\subsection*{Перешифрование PIN-блока (PIN translation)}

Когда PIN-блок передаётся между зонами безопасности (например,
от~эквайера к~карточной сети и~далее к~эмитенту), каждый участник
использует свой ключ шифрования. На~каждом стыке зон HSM выполняет
перешифрование PIN-блока: расшифровывает одним ключом и~тут же
зашифровывает другим, не~выпуская PIN в~память сервера. Операция
атомарна и~идёт исключительно внутри HSM. Поэтому любой узел,
обрабатывающий PIN, обязан иметь HSM.

\section{Церемония ключа (Key Ceremony)}\label{sec:crypto-key-ceremony}

Мастер-ключ~--- самый ценный секрет в~платёжной инфраструктуре,
его компрометация означает компрометацию всех производных ключей,
всех карт, всех транзакций. Процедуру генерации, загрузки и~хранения
мастер-ключа называют Key Ceremony~--- церемония генерации, разбиения
и~загрузки ключа с~dual control (двойной контроль: любая критическая операция требует одновременного участия двух операторов)
и~split knowledge (раздельное знание: ни один отдельный участник не~владеет полным ключом). Правило
жёсткое: ни один человек не~должен увидеть ключ целиком~\cite{pci-pin-security}.

Процедура Key Ceremony канонически состоит из четырёх шагов:

\begin{enumerate}
  \item HSM генерирует ключ и~делит его на~компоненты, ни один из~которых
    сам по~себе не~позволяет восстановить исходный секрет.
  \item Компоненты записывают на~отдельные носители и~передают
    разным хранителям.
  \item Хранители подтверждают получение и~хранят свои части отдельно.
  \item Для загрузки ключа компоненты собирают только внутри защищённого
    периметра HSM.
\end{enumerate}

Нарушение dual control или split knowledge на любом из шагов~--- инцидент
безопасности, требующий перегенерации ключа.

Для одного мастер-ключа с~тремя компонентами и~двумя HSM (primary плюс backup)
церемония занимает 4--8~часов: подготовка комнаты, генерация и~разбиение,
раздельная передача компонентов хранителям, загрузка в~резервный HSM
и~верификация, документирование и~опечатывание сейфов~\cite{pci-pin-security}.

Split knowledge защищает от~инсайдера: даже если один из
хранителей скомпрометирован или действует под принуждением,
у~него нет полного ключа. Dual control добавляет второй барьер:
ни одно действие с~мастер-ключом не~завершается без участия
второго человека. Атака на~мастер-ключ требует сговора
нескольких людей; одиночному злоумышленнику она недоступна.

\subsection*{Иерархия церемоний и~доверия}

Церемоний несколько, и~они образуют иерархические уровни. Чем
ближе ключ к~корню доверия, тем дороже его обращение и~тем строже
процедура; чем ниже~--- тем больше операций и~тем сильнее
автоматизация.

\begin{itemize}
  \item \textbf{Корневой уровень.} Local Master Key (LMK) внутри
    HSM эквайера или эмитента и~корневой ключ PKI (Public Key Infrastructure, инфраструктура открытых ключей), которым потом
    подписываются сертификаты сервера RKI (Remote Key Injection, удалённая загрузка ключей; разбирается ниже в~разделе \enquote{Загрузка ключей в~терминалы}), генерируются и~грузятся
    очной церемонией с~dual control и~split knowledge. То же относится
    к~корневому ключу CA (Certification Authority, удостоверяющий центр) эмитента, под которым выпускаются
    сертификаты карт EMV. Эти ключи живут долго (годы), и~их перевыпуск~---
    отдельный проект.
  \item \textbf{Межорганизационный уровень.} Zone Master Key (ZMK)
    или Key Encryption Key (KEK), которыми банк, процессор и~карточная
    сеть обмениваются ключами между собой, передаются компонентами
    под подписанные акты~--- тоже в~ходе формальной церемонии,
    хотя без~присутствия на~одной площадке (компоненты курьерятся
    в~опечатанных конвертах).
  \item \textbf{Производный уровень.} Терминальные ключи (TPK,
    Terminal PIN Key, для шифрования PIN-блока; TMK, Terminal Master Key, как локальный ключ хранения), сессионные ключи DUKPT и~ключи
    отдельных карт~--- производные. Они выпускаются автоматически
    HSM из~ключей верхнего уровня и~до участников доходят
    по~защищённым каналам без отдельной церемонии для каждого
    устройства.
\end{itemize}

\highlight{RKI переносит церемонию на~один уровень
выше}: вместо ручной загрузки ключа в~каждый терминал церемонией
один раз обслуживается корневой ключ сервера RKI, а~из него уже
автоматически производятся ключи отдельных терминалов.

\subsection*{Загрузка ключей в~терминалы (key injection)}

Ключи попадают в POS-терминалы одним из двух способов.

\begin{itemize}
  \item \textbf{Ручная загрузка (manual key injection)}~--- терминал
    привозят в~защищённую зону, где оператор загружает начальный
    ключ через специальное устройство. Процедура трудоёмка
    и~плохо масштабируется, но по-прежнему применяется для
    первоначальной инициализации.
  \item \textbf{Удалённая загрузка (Remote Key Injection, RKI)}~---
    ключ загружают по~защищённому каналу через специальный сервер.
    Терминал должен пройти взаимную аутентификацию и~принять
    только тот ключ, который пришёл по~доверенному каналу. RKI
    упрощает логистику, особенно для больших терминальных
    сетей. В~терминологии PCI PIN это удалённое распределение
    и~установление ключей с~использованием асимметрических
    техник из~Annex A~--- первичная загрузка ключей
    с~корневым доверенным сертификатом
    PKI~\cite{pci-pin-security}. Корневой ключ
    самого сервера RKI при этом грузится так же, как LMK,~---
    очной церемонией.
\end{itemize}

\section{ГОСТ в~платёжном приложении Мир}\label{sec:crypto-gost}

Российские требования добавляют жёсткое ограничение: процессинг,
умеющий только RSA, 3DES и~AES, к~полноценной работе с~картами
Мир не~готов. Требования к~защите информации в~платёжной системе
задают действующие нормы Банка России; локальная криптография
строится вокруг алгоритмов ГОСТ, точная регуляторная
раскладка разобрана в~разделе~\ref{sec:regulation-cbr-provisions}~\cite{cbr-821p,cbr-851p}.
Базовый набор криптографических ГОСТ системы Мир:

\begin{itemize}
  \item \textbf{ГОСТ~34.10--2018}~--- алгоритм электронной
    подписи на~эллиптических кривых. В~платежах он выполняет
    ту же роль, что и~другие схемы цифровой подписи
    в~международных системах, и~доказывает, что данные не~были
    подменены~\cite{gost-34-10}.
  \item \textbf{ГОСТ~34.11--2018}, он же Стрибог,~--- функция хеширования,
    используемая при построении подписи и~в~ряде криптографических
    протоколов~\cite{gost-34-11}.
  \item \textbf{ГОСТ~34.12--2018}~--- блочные шифры Кузнечик
    и~Магма, применяемые для шифрования данных и~вычисления имитовставки (имитовставка~--- эквивалент MAC в~отечественной терминологии)~\cite{gost-34-12}.
    В~российской криптографии Магма играет роль \emph{legacy}-шифра
    (наследник ГОСТ~28147--89), а~Кузнечик~--- \emph{current}-выбор для
    новых протоколов; пара симметрична международной разводке
    \enquote{3DES legacy / AES current}.
  \item \textbf{ГОСТ~34.13--2018}~--- режимы работы блочных
    шифров (простая замена, гаммирование, гаммирование с~обратной
    связью, выработка имитовставки и~другие), на~которых
    строятся протоколы шифрования и~имитозащиты~\cite{gost-34-13}.
\end{itemize}
Исторически основным симметричным шифром был ГОСТ~28147--89;
сейчас этот алгоритм входит в~ГОСТ~34.12--2018 под названием
\enquote{Магма}, а~ссылки на~28147--89 встречаются в~унаследованных (\textit{legacy}) конфигурациях
и~в~ряде межведомственных документов как опорный
алгоритм~\cite{gost-28147}.
Для процессинга это означает двойной стек (\textit{dual-stack})~--- HSM
одновременно поддерживает ГОСТ и~международные алгоритмы. На~кобейджинговых
картах Мир--UnionPay внутрироссийская транзакция идёт по~алгоритмам
ГОСТ, операция за~рубежом~--- по~международным. Процессинг эмитента
должен обрабатывать оба типа криптограмм, HSM~--- хранить и~использовать
оба набора ключей~\cite{gost-34-10,gost-34-12}.

Внедрение поддержки ГОСТ начинается с~HSM. Не все модели поддерживают
алгоритмы ГОСТ штатно: ряд западных HSM требует дополнительного
модуля или прошивки. После обновления HSM нужна отдельная церемония
ключа для набора ГОСТ, загрузка мастер-ключей ГОСТ и~настройка
деривации сессионных ключей по~схеме НСПК. Тестовая среда тоже
должна поддерживать оба набора~--- без этого кобейджинговые сценарии
не~воспроизвести.

\begin{figure}[tbp]
  \centering
  \includefiguresvg[width=.7\textwidth]{assets/figures/ch10-cryptography/gost-vs-international}
  \caption{Криптография ГОСТ и~международные алгоритмы по~ролям:
  подпись, хеш, блочный шифр, режим.}\label{fig:gost-vs-international}
\end{figure}

Операционные детали самой системы Мир, внутренняя маршрутизация
и кобейджинговые сценарии разобраны в~главе~\ref{ch:mir}.
